% answers_short.tex —— 2015 年·广东卷·理科 · 简短答案速查
% 由 tools/make_short_answers.py 自 answers.tex 的 \gksolhead/\gkanswer 提取，
% 并与 2015/2015_数学真题_广东理_答案要点.md 逐题比对后生成，勿手改（改源 MD 后重跑）。

\begin{gknotice}
\textbf{使用说明}：本节只给最终答案便于快速核对，完整推导见解析版。答案为编写者独立推导，\textbf{非官方标准答案}。
\end{gknotice}

\gkpart{选择题}
\begin{center}\begin{tabular}{ccccccccc}
\toprule
题号 & 1 & 2 & 3 & 4 & 5 & 6 & 7 & 8 \\
\midrule
答案 & D & A & D & B & A & B & C & B \\
\bottomrule
\end{tabular}\end{center}

\gkpart{填空题}
\begin{itemize}[leftmargin=2.6em]
  \item[\textbf{9.}] \(6\)
  \item[\textbf{10.}] \(10\)
  \item[\textbf{11.}] \(1\)
  \item[\textbf{12.}] \(1560\)
  \item[\textbf{13.}] \(\dfrac{1}{3}\)
  \item[\textbf{14.}] \(\dfrac{5\sqrt{2}}{2}\)
  \item[\textbf{15.}] \(8\)
\end{itemize}

\gkpart{解答题}
\begin{itemize}[leftmargin=2.6em]
  \item[\textbf{16.}] (1) \(\tan x=1\)（此时 \(x=\dfrac{\pi}{4}\)）；(2) \(x=\dfrac{5\pi}{12}\)
  \item[\textbf{17.}] (1) 样本年龄 \(44,40,36,43,36,37,44,43,37\)（编号 \(2,6,10,14,18,22,26,30,34\)）；(2) \(\bar{x}=40\)，\(s^2=\dfrac{100}{9}\)；(3) \(23\) 人，占 \(63.89\%\)
  \item[\textbf{18.}] (1) 成立（\(PE\perp\) 平面 \(ABCD\)，\(FG\subset\) 平面 \(ABCD\)，故 \(PE\perp FG\)）；(2) \(\dfrac{\sqrt{7}}{3}\)；(3) \(\dfrac{9\sqrt{5}}{25}\)
  \item[\textbf{19.}] (1) 单调递增区间为 \((-\infty,+\infty)\)，无递减区间；(2) 成立（\(f\) 在 \(\mathbf{R}\) 上严格递增，且 \(x\to-\infty\) 时 \(f(x)\to-a<0\)、\(x\to+\infty\) 时 \(f(x)\to+\infty\)）；(3) 成立（由 \(\mathrm{e}^m\geqslant m+1\) 得 \(a-\dfrac{2}{\mathrm{e}}=\mathrm{e}^m(m+1)^2\geqslant(m+1)^3\)，再开立方）
  \item[\textbf{20.}] (1) 圆心 \(C_1(3,0)\)；(2) 轨迹 \(C\)：\(x^2+y^2-3x=0\)，且 \(\dfrac{5}{3}<x\leqslant3\)；(3) 存在，\(k\in\left[-\dfrac{2\sqrt{5}}{7},\dfrac{2\sqrt{5}}{7}\right]\cup\left\{-\dfrac{3}{4},\dfrac{3}{4}\right\}\)
  \item[\textbf{21.}] (1) \(a_3=\dfrac{1}{4}\)；(2) \(T_n=2-2^{1-n}\)；(3) 成立（\(S_n=H_n\left(2-2^{1-n}\right)<2H_n<2+2\ln n\)）
\end{itemize}
